\documentclass[12pt,a4paper]{article}
\usepackage{expediente}
\begin{document}
% =============================================================================
% FICHA DE DATOS DEL EXPEDIENTE --- actualizado 20 ago 2026
% =============================================================================

% --- VENDEDOR / TITULAR REGISTRAL -------------------------------------------
\newcommand{\VendedorNombres}{RICHAR DARWIN}
\newcommand{\VendedorApellidos}{CUTIRE CONDORI}
\newcommand{\VendedorNombreCompleto}{RICHAR DARWIN CUTIRE CONDORI}
\newcommand{\VendedorDNI}{42610690}
\newcommand{\VendedorEstadoCivil}{soltero, sin unión de hecho vigente}
\newcommand{\VendedorOcupacion}{\faltante{[ocupación de Richar]}}
\newcommand{\VendedorDomicilio}{\faltante{[domicilio actual de Richar]}}
\newcommand{\VendedorTelefono}{\faltante{[teléfono de Richar]}}
\newcommand{\VendedorCorreo}{\faltante{[correo]}}
\newcommand{\VendedorNacionalidad}{peruano}
\newcommand{\VendedorFechaNac}{\faltante{[fecha de nacimiento]}}
\newcommand{\VendedorDNIEmision}{\faltante{[fecha de emisión del DNI de Richar --- la pide SUNARP]}}

% --- CONVIVIENTE: NO CORRESPONDE -------------------------------------------
\newcommand{\ConvivienteNombre}{NO CORRESPONDE}
\newcommand{\ConvivienteDNI}{NO CORRESPONDE}
\newcommand{\ConvivienteDesde}{NO CORRESPONDE}
\newcommand{\ConvivienteInscrita}{no existe unión de hecho vigente}
\newcommand{\ConvivienteDomicilio}{NO CORRESPONDE}

% --- COMPRADORES (copropiedad en partes iguales) ---------------------------
\newcommand{\CompradorUnoNombre}{NICO ALVARO CUTIRE ARCE}
\newcommand{\CompradorUnoDNI}{48325017}
\newcommand{\CompradorUnoEmision}{01 de enero de 2024}
\newcommand{\CompradorUnoEstadoCivil}{soltero}
\newcommand{\CompradorUnoOcupacion}{ingeniero informático}
\newcommand{\CompradorUnoDomicilio}{Jr.\ Wiracocha, Sicuani, provincia de Canchis, Cusco, Perú}
\newcommand{\CompradorUnoTelefono}{984715108}

\newcommand{\CompradorDosNombre}{LUCY CUTIRE ARCE}
\newcommand{\CompradorDosDNI}{47478852}
\newcommand{\CompradorDosEmision}{09 de enero de 2026}
\newcommand{\CompradorDosEstadoCivil}{soltera}
\newcommand{\CompradorDosOcupacion}{contador}
\newcommand{\CompradorDosDomicilio}{Comunidad Hanocca, distrito de Layo, provincia de Canas, Cusco, Perú}
\newcommand{\CompradorDosTelefono}{953440875}

\newcommand{\CompradorNombres}{NICO ALVARO y LUCY}
\newcommand{\CompradorApellidos}{CUTIRE ARCE}
\newcommand{\CompradorNombreCompleto}{NICO ALVARO CUTIRE ARCE y LUCY CUTIRE ARCE}
\newcommand{\CompradorDNI}{48325017 y 47478852}
\newcommand{\CompradorEstadoCivil}{solteros}
\newcommand{\CompradorConyuge}{NO CORRESPONDE}
\newcommand{\CompradorConyugeDNI}{NO CORRESPONDE}
\newcommand{\CompradorOcupacion}{ingeniero informático y contador, respectivamente}
\newcommand{\CompradorDomicilio}{Sicuani (Canchis) y Layo (Canas), Cusco}
\newcommand{\CompradorTelefono}{984715108 / 953440875}
\newcommand{\CompradorCorreo}{\faltante{[correo de contacto]}}
\newcommand{\CompradorNacionalidad}{peruanos}
\newcommand{\PorcentajeCopropiedad}{50\,\% cada uno}

% --- INTERMEDIARIO / OCUPANTE DEL REMANENTE --------------------------------
\newcommand{\IntermediarioNombre}{LUCIO GIL CUTIRE PARI}
\newcommand{\IntermediarioDNI}{42467850}
\newcommand{\IntermediarioEstadoCivil}{\faltante{[estado civil de Lucio]}}
\newcommand{\IntermediarioOcupacion}{\faltante{[ocupación]}}
\newcommand{\IntermediarioDomicilio}{\faltante{[domicilio de Lucio]}}
\newcommand{\IntermediarioTelefono}{\faltante{[teléfono de Lucio]}}
\newcommand{\IntermediarioParentesco}{familiar (apellido Cutire)}

% --- PREDIO MATRIZ ----------------------------------------------------------
\newcommand{\PartidaMatriz}{\faltante{[N.º de partida electrónica SUNARP --- pendiente]}}
\newcommand{\OficinaRegistral}{Oficina Registral de Puerto Maldonado}
\newcommand{\AreaMatriz}{\faltante{[área total inscrita; las partes refieren del orden de 100 ha]}}
\newcommand{\UnidadCatastral}{\faltante{[código catastral / UC, si existe]}}
\newcommand{\NombreFundo}{\faltante{[nombre del fundo, si tiene]}}
\newcommand{\KmCarretera}{Interoceánica Sur, al este del Puente Jayave (el puente queda al oeste del predio, sobre el río; no debajo del lote)}

% --- PREDIO A INDEPENDIZAR --------------------------------------------------
\newcommand{\AreaIndependizarHa}{10{,}00}
\newcommand{\AreaIndependizarMetros}{100 000{,}00}
\newcommand{\PerimetroM}{2 200{,}58}
\newcommand{\FrenteM}{100{,}29}
\newcommand{\FondoM}{1 000{,}00}

\newcommand{\PUnoLat}{$-12{,}912452$}
\newcommand{\PUnoLon}{$-70{,}170058$}
\newcommand{\PDosLat}{$-12{,}912497$}
\newcommand{\PDosLon}{$-70{,}170981$}
\newcommand{\PTresLat}{$-12{,}903469$}
\newcommand{\PTresLon}{$-70{,}171438$}
\newcommand{\PCuatroLat}{$-12{,}903424$}
\newcommand{\PCuatroLon}{$-70{,}170515$}

\newcommand{\PUnoE}{373 065{,}00}
\newcommand{\PUnoN}{8 572 256{,}00}
\newcommand{\PDosE}{372 965{,}00}
\newcommand{\PDosN}{8 572 251{,}00}
\newcommand{\PTresE}{372 911{,}00}
\newcommand{\PTresN}{8 573 249{,}00}
\newcommand{\PCuatroE}{373 011{,}00}
\newcommand{\PCuatroN}{8 573 255{,}00}

\newcommand{\FechaPago}{15 de agosto de 2018}
\newcommand{\MontoPago}{S/\ 20 000{,}00}
\newcommand{\MontoPagoLetras}{VEINTE MIL Y 00/100 SOLES}
\newcommand{\FormaPagoHistorica}{dinero en efectivo, sin medio de pago del sistema financiero}

% Nombres genéricos --- sustituir por los reales al ir a notaría
\newcommand{\NotarioNombre}{Dr.\ JUAN CARLOS RÍOS SALINAS (nombre genérico a sustituir)}
\newcommand{\NotariaCiudad}{Puerto Maldonado}
\newcommand{\ColegiaturaAbogado}{CAL Madre de Dios N.º 0000 (genérico)}
\newcommand{\AbogadoNombre}{Abog.\ ANA MARÍA QUISPE HUAMÁN (nombre genérico a sustituir)}
\newcommand{\IngenieroNombre}{Ing.\ PEDRO ANTONIO MAMANI CONDORI (nombre genérico a sustituir)}
\newcommand{\IngenieroCIP}{CIP 000000 (genérico)}
\newcommand{\Municipalidad}{Municipalidad Distrital de Inambari}
\newcommand{\Provincia}{Tambopata}
\newcommand{\Departamento}{Madre de Dios}
\newcommand{\Distrito}{Inambari}

\newcommand{\LugarOtorgamiento}{Puerto Maldonado}
\newcommand{\DiaOtorgamiento}{\faltante{[día]}}
\newcommand{\MesOtorgamiento}{\faltante{[mes]}}
\newcommand{\AnioOtorgamiento}{2026}

\newcommand{\TestigoUno}{\faltante{[testigo 1: nombres, DNI, domicilio]}}
\newcommand{\TestigoDos}{\faltante{[testigo 2: nombres, DNI, domicilio]}}

\InicioDocumento{Acta de linderación, amojonamiento y entrega de posesión del lote de 10,00 ha}{DOC-18}

\noindent En \LugarOtorgamiento, siendo las \faltante{[hora]} horas del \DiaOtorgamiento\ de \MesOtorgamiento\ de \AnioOtorgamiento, en el predio rústico sito al norte de la Interoceánica Sur, al este del Puente Jayave, distrito de \Distrito, se reunieron:

\begin{itemize}[leftmargin=1.1cm]
\item \textbf{EL TITULAR:} \VendedorNombreCompleto, DNI \VendedorDNI.
\item \textbf{EL COMPRADOR:} \CompradorNombreCompleto, DNI \CompradorDNI.
\item \textbf{EL OCUPANTE:} \IntermediarioNombre, DNI \IntermediarioDNI.
\item \textbf{EL PROFESIONAL:} \IngenieroNombre, CIP \IngenieroCIP.
\item \textbf{TESTIGOS:} \TestigoUno; y \TestigoDos.
\end{itemize}

\Clausula{Primero --- Recorrido}
Los comparecientes recorrieron el polígono de trabajo P1--P2--P3--P4, identificando el frente sobre la carretera y el fondo hacia el interior. Se constató / se dejará constancia de:

\vspace{0.2em}
\noindent $\square$ coincidencia con el norte de la vía \quad $\square$ necesidad de invertir al sur \quad $\square$ ajuste de hitos en campo (croquis anexo).

\Clausula{Segundo --- Amojonamiento}
Se colocaron / se acordó colocar hitos en:

{\small\renewcommand{\arraystretch}{1.2}
\noindent\begin{tabularx}{\textwidth}{c X c}
\toprule
Vértice & Descripción en campo & Hito (sí/no) \\
\midrule
P1 & Frente este, borde vía & \linea[2.4cm] \\
P2 & Frente oeste, borde vía & \linea[2.4cm] \\
P3 & Fondo oeste & \linea[2.4cm] \\
P4 & Fondo este & \linea[2.4cm] \\
\bottomrule
\end{tabularx}}

\Clausula{Tercero --- Entrega de posesión}
EL OCUPANTE \textbf{entrega} a EL COMPRADOR la posesión inmediata, pacífica y pública del lote de 10,00 ha, a título de futuro propietario / propietario según la escritura que se otorgue. EL COMPRADOR \textbf{recibe} dicha posesión. EL TITULAR presta su conformidad.

\vspace{0.25em}
\noindent Fecha efectiva de entrega: \faltante{[hoy / otra fecha]}.\\
Bienes, semovientes, campamentos o plantaciones que se retiran o se quedan: \faltante{[detalle]}.

\Clausula{Cuarto --- Oposición}
Durante el acto \textbf{no} se formuló oposición de tercero / se formuló la siguiente: \faltante{[ninguna / detalle]}.

\Clausula{Quinto --- Fotografías y coordenadas}
Se tomaron fotografías de los hitos y de las partes, que se anexan. El profesional registró coordenadas de campo que \textbf{prevalecen} sobre el DOC-08 si hubiere diferencia, debiendo actualizarse memorias y minutas.

\Clausula{Sexto --- Lectura}
Leída el acta, la suscriben en señal de conformidad, junto con su huella dactilar.

\LugarFecha
\TresFirmas{TITULAR}{\VendedorNombreCompleto\\DNI \VendedorDNI}{COMPRADOR}{\CompradorNombreCompleto\\DNI \CompradorDNI}{OCUPANTE}{\IntermediarioNombre\\DNI \IntermediarioDNI}

\vspace{0.55em}
\noindent
\BloqueFirmaSimple{PROFESIONAL}{\IngenieroNombre\\CIP \IngenieroCIP}
\hfill
\BloqueFirmaSimple{TESTIGO 1}{\TestigoUno}

\vspace{0.45em}
\noindent
\BloqueFirmaSimple{TESTIGO 2}{\TestigoDos}

\vspace{0.7em}
\noindent{\small\textit{Croquis de campo / fotos}}\\
\noindent\fbox{\parbox[c][3.3cm][c]{0.96\textwidth}{\centering\color{black!40}espacio para croquis, fotos de hitos y de las partes en el predio}}

\end{document}
